\documentclass[11pt]{article}
\usepackage[margin=1in]{geometry}
\usepackage{booktabs,amsmath,amssymb,graphicx,hyperref,natbib}
\usepackage[T1]{fontenc}
\newcommand{\Hcn}{\ensuremath{H_{\mathrm{cn}}}}
\title{\bfseries How Much Ranking a Benchmark Can Support Is Predictable Before It Is Run}
\author{Ilpo V\"a\"at\"ainen\thanks{Independent researcher, Helsinki. \texttt{ipezygj2@gmail.com}. Code and results: \url{https://github.com/ipezygj/swebench-rank-audit}.}}
\date{August 2026}
\begin{document}
\maketitle

\begin{abstract}
A leaderboard reports an ordering, but how much of that ordering the
evidence actually supports is a property of the measuring instrument
rather than of the systems on it. We show that two such properties are
predictable in closed form from four numbers available before any system
is run: the number of entrants $J$, the number of items $n$, the spread of
scores $\tau$, and the paired difference SD $\sigma_p$.
The share of pairs the evidence establishes follows
$\bar\Phi(1/\mathrm{SNR})$ with $\mathrm{SNR}=\tau\sqrt{2n}/(c\,\sigma_p)$,
to a mean absolute error of 5.1 points across 9 leaderboards in
five fields (3.3 points using a robust estimate of $\tau$), and
the entropy of the orderings the evidence permits is reproduced by a
Gaussian field with the same four numbers and nothing else of the real
board, within five points on 6 of 9 boards.
A tenth board was held out until the thresholds were committed and passes
both. Neither law is fitted: the only constant is the critical value the
multiplicity procedure itself uses, and when that procedure was replaced
after its coverage was found wanting -- moving the realised critical value
from 3.14 to 8.45 across boards -- the predictions followed it.
The laws are neutral about whether a benchmark is good: they say what it
can decide, and for two of the boards here the answer is almost nothing
while for another it is everything.
\end{abstract}

\section{Introduction}

A leaderboard prints an order. Some of that order is evidence and some of
it is sampling noise on a finite item set, and the two are not
distinguishable by looking at the table. This is not a new observation, and
the tools for acting on it are not new either.

\citet{card2020} measured the statistical power of NLP experiments and
found that most attempted comparisons to state of the art on the GLUE tasks
are underpowered, and that a 2\,000-sentence machine-translation test set
has about 75\,\% power for a one-point BLEU difference. \citet{mogstad2024}
give simultaneous confidence sets for ranks, and two 2026 papers carry that
construction to model leaderboards \citep{rankintervals}.
\citet{uncertaintyranking} quantify rank uncertainty by counting the linear
extensions of the order induced by interval estimates.
\citet{saturation} define an uncertainty-aware saturation index from
leaderboard data and observe that discriminative power is lost once the top
models sit within a couple of points.

Every one of these measures a leaderboard that already exists. Power
analysis needs the observed variance; rank intervals need the observed
scores; the saturation index takes the highest observed score as its
ceiling. That is the right thing to do when auditing a published table, and
it leaves one question unanswered: \emph{how much ranking will a benchmark
be able to support}, asked by someone deciding how many items to write
before any system has been run on any of them.

This paper answers that question with two closed forms in four quantities:
the number of entrants $J$, the number of items $n$, the spread of the
field's scores $\tau$, and the difference SD of a typical pair $\sigma_p$.
Three of the four are design choices or planning estimates; the fourth,
$\sigma_p$, is a property of the item set that a pilot on two systems
measures. Nothing is fitted to leaderboards: the only constant in the
formula is the critical value the multiplicity procedure itself returns.

\paragraph{Contributions.}
\begin{enumerate}
\item A closed form for the share of pairs a leaderboard's evidence
  establishes, $\bar\Phi(1/\mathrm{SNR})$ with
  $\mathrm{SNR} = \tau\sqrt{2n}/(c\,\sigma_p)$, holding to a mean absolute
  error of 5.1 points (3.3 robust) across 9 leaderboards.
\item The finding that the entropy of the orderings the evidence permits is
  reproduced by a Gaussian field with the same four numbers and nothing else
  of the real board, on 6 of 9 boards.
\item Validation across five fields -- code, embeddings, competition
  mathematics, protein fitness, tabular prediction -- and on a tenth board
  held out until every threshold had been committed to version control.
\item Evidence that the laws are about resolution rather than about an
  estimator: the multiplicity procedure underneath was found to undercover,
  was replaced, the realised critical value moved from 3.14 to 8.45 across
  boards, and the predictions followed (\S\ref{sec:robust}).
\end{enumerate}

\paragraph{What is not claimed.} That leaderboards are bad; the laws are
neutral and say of one board here that it resolves its top at $t = 9.89$.
That benchmark comparisons are underpowered, which is \citet{card2020}.
That we contribute a method for rank inference, which is
\citet{mogstad2024}. And the laws do not predict the frontier: they
reproduce aggregates and miss, by a factor of six on one board, the number
of systems that could be first.

\section{What is being predicted}

Fix a leaderboard: $J$ systems, each scored on the same $n$ items, so that
every pair can be compared on paired differences. Write $x_{ji}$ for system
$j$'s outcome on item $i$, $s_j$ for its mean, and for a pair $(j,k)$ let
$d_i = x_{ji} - x_{ki}$ with SD $\sigma_{jk}$. Two summaries of what the
evidence supports are then well defined.

\paragraph{The established share.} The fraction of ordered pairs $(j,k)$
for which the evidence places $j$ above $k$ at a simultaneous level. This
is what a reader is entitled to read off the printed order.

\paragraph{The ordering entropy.} $\log_2$ of the number of total orders
consistent with the established partial order, normalised by $\log_2 J!$.
Zero means the evidence fixes the ranking; one means it fixes nothing.
This quantity is not new -- \citet{uncertaintyranking} measure rank
uncertainty by counting linear extensions of the order induced by interval
estimates -- and counting them exactly is \#P-complete, so we use Knuth's
estimator.

Neither summary is the contribution. The contribution is that both are
predictable from $(J, n, \tau, \sigma_p)$ before the board exists.

\section{Law 1: the established share}

A simultaneous procedure at level $\alpha$ returns a critical value $c$ and
declares $j$ above $k$ when
\begin{equation}
|s_j - s_k| \;>\; c\,\frac{\sigma_{jk}}{\sqrt{n}} .
\label{eq:sep}
\end{equation}
Treat the pair as drawn from the field. The difference of two scores drawn
from a spread of SD $\tau$ has SD $\tau\sqrt{2}$, so writing $\sigma_p$ for
the difference SD of a typical pair, the probability that a random pair
clears \eqref{eq:sep} is
\begin{equation}
\Pr\!\left(|s_j - s_k| > \frac{c\,\sigma_p}{\sqrt{n}}\right)
\;=\; 2\,\bar\Phi\!\left(\frac{c\,\sigma_p}{\sqrt{2n}\,\tau}\right)
\;=\; 2\,\bar\Phi(1/\mathrm{SNR}),
\qquad
\mathrm{SNR} \;=\; \frac{\tau\sqrt{2n}}{c\,\sigma_p} .
\label{eq:law1}
\end{equation}

A separated \emph{unordered} pair contributes exactly one entry to the
$J(J-1)$ ordered pairs, so the established share as reported here -- the
density of the beats matrix -- is half of \eqref{eq:law1}, namely
$\bar\Phi(1/\mathrm{SNR})$. We labour the factor of two because the first
version of this work predicted the unordered share, compared it against the
ordered one, and failed its own Gaussian self-check by exactly $2\times$;
the definitions were reconciled, not the data.

$\mathrm{SNR}$ has a reading: it is the spread of the field measured in
units of one pair's simultaneous resolution. It grows as $\sqrt{n}$, so
quadrupling the item set doubles it, and it falls as $c$ grows, which is
how the price of comparing many systems at once enters.

\paragraph{What the derivation assumes.} Three things, and they are worth
naming because the boards that fail the law fail exactly one of them.

\begin{description}
\item[A1: one $\sigma_p$ for all pairs.] False in detail. On five dated
  boards the pair-specific difference SD varies by 6 to 47\,\% between pairs,
  and substituting the board-wide figure into a power calculation misstates
  the items needed by between $1.0\times$ and $45\times$ for individual
  pairs. It survives here because the law is an aggregate: integrating
  \eqref{eq:law1} over the observed distribution of $\sigma_{jk}$ instead of
  using its median changes the mean absolute error from 5.1 to 4.3 points.
  The heterogeneity averages out of the aggregate and matters entirely for
  the individual claim.
\item[A2: the score differences are Gaussian.] This is the load-bearing
  assumption and the one that breaks. It enters only through $\bar\Phi$,
  and a field with a long lower tail has a $\tau$ inflated by systems far
  from the pack, so \eqref{eq:law1} over-predicts. Both TabArena boards
  behave this way, and replacing $\tau$ by an IQR-based estimate -- which
  ignores the tails -- moves their errors from $+17.8$ and $+17.2$ points
  to $-8.5$ and $-8.1$.
\item[A3: $c$ is a scalar.] Exactly true for a single-step procedure and
  approximately true for a step-down one, where we take $c$ to be the
  realised critical value at the final step.
\end{description}

\paragraph{What the derivation does \emph{not} assume.} Independence between
systems. $\sigma_p$ is the SD of the \emph{paired difference} series, which
already contains whatever correlation the systems have; two systems that
fail the same items have a small $\sigma_{jk}$ and separate more easily at
the same score gap. This is why the law needs no term for the dependence
between entrants, and it is consistent with a separate finding of ours that
a twin carrying the real residual correlation structure is
indistinguishable from an independent one on these aggregates.

\begin{table}[t]\centering\small
\begin{tabular}{l rr r rr rr}
\toprule
 & $J$ & $n$ & SNR & observed & predicted & err & err (robust) \\\\
\midrule
SWE-bench Verified & 134 & 500 & 3.00 & 37.4\,\% & 36.9\,\% & $-0.4$ & $+0.6$ \\
MTEB English v2 & 181 & 41 & 1.96 & 28.5\,\% & 30.5\,\% & $+1.9$ & $-3.5$ \\
HELM classic & 90 & 10 & 0.58 & 6.6\,\% & 4.3\,\% & $-2.2$ & $+1.2$ \\
ProteinGym DMS & 96 & 217 & 2.70 & 34.4\,\% & 35.5\,\% & $+1.1$ & $-4.4$ \\
TabArena 16 models & 16 & 51 & 3.12 & 19.6\,\% & 37.4\,\% & $+17.8$ & $-8.5$ \\
TabArena 45 variants & 45 & 51 & 2.99 & 19.7\,\% & 36.9\,\% & $+17.2$ & $-8.1$ \\
CASP14 & 101 & 42 & 1.51 & 21.0\,\% & 25.4\,\% & $+4.4$ & $-0.8$ \\
LiveBench & 152 & 200 & 1.69 & 28.8\,\% & 27.7\,\% & $-1.2$ & $+1.3$ \\
MathArena 2025 & 35 & 183 & 1.94 & 30.3\,\% & 30.3\,\% & $+0.0$ & $-1.0$ \\
\bottomrule
\end{tabular}
\caption{Law 1 across 9 leaderboards in five fields. The predicted column uses the SD of scores as $\tau$; the robust error column uses the IQR rescaled to a Gaussian SD, which ignores the tails of the field. Mean absolute error 5.1 points, or 3.3 points robust. No constant in the formula is fitted.}
\label{tab:law1}
\end{table}

\section{Law 2: the ordering entropy}

Law 1 is a formula because the established share counts pairs. Entropy
does not: it depends on \emph{which} pairs are established, since a partial
order concentrated among the weakest systems constrains fewer total orders
than the same number of relations spread through the field. Counting linear
extensions is \#P-complete, and we know of no closed form. The claim here is
therefore of a different and weaker kind, and it is stated as such:

\begin{quote}
A Gaussian field with the same $J$, $n$, $\tau$ and $\sigma_p$ as a real
leaderboard, and nothing else of it, reproduces its ordering entropy.
\end{quote}

The twin is constructed by inverting the noise: the observed spread $\tau$
already contains measurement error, so the latent spread is
$\sqrt{\max(\tau^2 - \sigma_{\mathrm{item}}^2/n,\,0)}$ with
$\sigma_{\mathrm{item}} = \sigma_p/\sqrt{2}$. Abilities are drawn from it,
item-level noise is added, and the identical machinery is run on the result.
The twin has no skew, no clusters and no outliers; if the real board's
entropy needed any of those, the twin would miss.

Two self-checks guard the construction: a twin of a Gaussian field must
reproduce its own entropy within three points on two independent seeds, and
the twin must reproduce its target's $\tau$ and $\sigma_p$ within 10\,\%.

\begin{table}[t]\centering\small
\begin{tabular}{l rr rrr}
\toprule
 & $J$ & $n$ & real & Gaussian twin & difference \\
\midrule
SWE-bench Verified & 134 & 500 & 54.6\,\% & 53.0\,\% & $+1.6$ \\
MTEB English v2 & 181 & 41 & 61.2\,\% & 61.8\,\% & $-0.6$ \\
HELM classic & 90 & 10 & 85.9\,\% & 82.8\,\% & $+3.1$ \\
ProteinGym DMS & 96 & 217 & 53.4\,\% & 53.2\,\% & $+0.2$ \\
TabArena 16 models & 16 & 51 & 63.9\,\% & 35.7\,\% & $+28.2$ \\
TabArena 45 variants & 45 & 51 & 67.2\,\% & 44.4\,\% & $+22.8$ \\
CASP14 & 101 & 42 & 75.0\,\% & 64.6\,\% & $+10.5$ \\
LiveBench & 152 & 200 & 65.3\,\% & 65.3\,\% & $+0.0$ \\
MathArena 2025 & 35 & 183 & 57.7\,\% & 54.3\,\% & $+3.4$ \\
\bottomrule
\end{tabular}
\caption{Law 2. Ordering entropy as a share of $\log_2 J!$, real against a twin that knows only $J$, $n$, $\tau$ and $\sigma_p$. Within five points on 6 of 9 boards.}
\label{tab:law2}
\end{table}

\section{Results}

\paragraph{Law 1.} Table~\ref{tab:law1} gives the nine boards. The mean
absolute error is 5.1 points, or 3.3 using a robust estimate of
$\tau$, against a quantity that ranges over the table from
37.4\,\% down to 6.6\,\%.
6 of the 9 sit within 2.5 points. The prediction is not a fit:
there is no free parameter in \eqref{eq:law1} to have absorbed the error.

\paragraph{Law 2.} Table~\ref{tab:law2} gives the same boards. The twin is
within five points on 6 of 9. Where it matches, a leaderboard's evidential
entropy is a function of four numbers and nothing else about the field --
which means a benchmark designer can compute it from a planned size and an
expected spread before a single system is run.

\paragraph{Both failures are the same failure.} The two TabArena boards
miss law 1 by about seventeen points and law 2 by twenty-three and
twenty-eight, and the direction was pre-registered before the run: the twin
has a symmetric field and spreads established pairs evenly, while the real
field concentrates them, and concentrated relations constrain fewer
orderings, so the \emph{real} entropy should come out higher. It does.

Decomposing the law-2 residual into a shape term -- what a twin gains when
it is given the real score shape instead of a Gaussian one -- and the
remainder separates the two cases cleanly:

\begin{table}[t]\centering\small
\begin{tabular}{l rrr rr}
\toprule
 & real & Gaussian & + real shape & SHAPE & residual \\
\midrule
SWE-bench Verified & 54.6\,\% & 53.0\,\% & 57.2\,\% & $+4.2$ & $-2.6$ \\
MTEB English v2 & 61.5\,\% & 61.7\,\% & 67.2\,\% & $+5.5$ & $-5.8$ \\
LiveBench & 66.0\,\% & 64.7\,\% & 66.2\,\% & $+1.5$ & $-0.2$ \\
TabArena 16 models & 63.6\,\% & 35.6\,\% & 75.1\,\% & $+39.5$ & $-11.5$ \\
TabArena 45 variants & 67.4\,\% & 44.6\,\% & 68.3\,\% & $+23.7$ & $-0.9$ \\
CASP14 & 74.5\,\% & 65.0\,\% & 73.7\,\% & $+8.7$ & $+0.8$ \\
\bottomrule
\end{tabular}
\caption{Where law 2's residual comes from. The SHAPE column is what the twin gains from being handed the real distribution of scores; the residual is what remains. On the boards the law fits, SHAPE is a few points and the residual cancels it; on TabArena's sixteen models SHAPE is $+39.5$.}
\label{tab:decomp}
\end{table}

\paragraph{It is the shape, not the noise profile.} A natural repair is to
give the twin each system's own item-noise level rather than one level for
the whole field. It does not help: that variant is within five points on
6/9 boards, the same count, and on the two TabArena boards the deviation
is not halved -- it moves from $+28.2$ to $+29.2$ on one and from $+22.9$ to
$+16.1$ on the other, against a pre-registered requirement that both be at
least halved. What the twin lacks is the asymmetry of the field, and giving
it a better noise model does not supply that.

\section{A board held out until the thresholds were committed}

Every threshold above was written into version control before the tenth
board was fetched. Its outcome:

\begin{table}[t]\centering\small
\begin{tabular}{l l l}
\toprule
& prediction & outcome \\
\midrule
\checkmark & kappa & over all pairs 1.00 +-0.05                   1.017 \\
\checkmark & kappa(\#1,\#2) & below 0.90                            0.758 \\
\checkmark & at & least 5 rank sets contain 1                     5 \\
\checkmark & law & 1 within 5 points                              observed 36.6 \%, law 35.6 \% \\
\checkmark & law & 2 within 5 points                              real 47.5 \%, twin 46.1 \% \\
$\times$ & pairing & dividend above 1.3                         1.00 (paired 10, independent 10) \\
$\times$ & kappa & split-half above 0.5 (expected to FAIL at n=28) r 0.41 \\
\bottomrule
\end{tabular}
\caption{The held-out board. The two failures were pre-registered as expected failures at that item count.}
\label{tab:tenth}
\end{table}

\section{Where the laws fail, and why that is informative}

\paragraph{Skewed fields.} Both TabArena boards miss law 1 by about
seventeen points and law 2 by twenty-three to twenty-eight, in the
direction pre-registered before the run: the Gaussian twin has a symmetric
field and spreads established pairs evenly, while the real field
concentrates them, and concentrated pairs constrain fewer orderings, so the
real entropy is higher. Two of nine.

\paragraph{The top.} The laws reproduce aggregates and miss the top. A
simulated board with SWE-bench Verified's shape reproduces its established
share and its entropy and then says three systems could be first where
nineteen can. Whatever governs the frontier is not in these four numbers.

\paragraph{$c$ is not free.} The law predicts resolution \emph{given} a
multiplicity procedure. That is a restriction, not a hidden parameter: $c$
is measured from the procedure, and \S\ref{sec:robust} is what happens when
the procedure changes.

\section{The laws survive replacing the estimator}
\label{sec:robust}

The rank sets underneath these numbers were originally built with a
multiplier bootstrap. Prompted by \citet{rankintervals}, who report that
bootstrap rank intervals fail under ties, we measured its simultaneous
coverage at the shapes actually used and found it failing on eight of
twelve boards -- 0.013 on the board with 90 systems and 10 items, 0.540 on
the one with 181 systems and 41 items, against a nominal 0.95 -- while Holm
on directional paired tests holds coverage on every shape. The construction
was replaced and the whole pipeline rerun.

This is the strongest test the laws have been put to, because the realised
critical value moved from 3.14 to 8.45 across boards and constructions. The
predictions followed: law 1's robust mean absolute error went from 3.4 to
3.3 points, and law 2 stayed within five points on the same six boards.
Of the 95 results files in the repository, 41 changed and 53 were
identical -- the ones whose measurements never touch rank sets.

\bibliographystyle{plainnat}
\begin{thebibliography}{9}
\bibitem[Card et al.(2020)]{card2020} D.~Card, P.~Henderson, U.~Khandelwal, R.~Jia, K.~Mahowald, D.~Jurafsky. With Little Power Comes Great Responsibility. \emph{EMNLP}, 2020.
\bibitem[Mogstad et al.(2024)]{mogstad2024} M.~Mogstad, J.~P.~Romano, A.~M.~Shaikh, D.~Wilhelm. Inference for Ranks with Applications to Mobility across Neighbourhoods and Academic Achievement across Countries. \emph{Review of Economic Studies}, 2024.
\bibitem[Rank Intervals(2026)]{rankintervals} Rank Intervals for Leaderboards: A Hierarchical Framework for Model Evaluation. arXiv:2606.08679, 2026.
\bibitem[Benchmark Saturation(2026)]{saturation} When AI Benchmarks Plateau: A Systematic Study of Benchmark Saturation. arXiv:2602.16763, 2026.
\bibitem[Uncertainty in Ranking(2021)]{uncertaintyranking} Uncertainty in Ranking. arXiv:2107.03459, 2021.
\end{thebibliography}
\end{document}
